\section{The auxiliary-unit extension graph does not close formally}
\label{sec:kt-polarized-extension-graph}

The three auxiliary-field units now have the precise
arithmetic consequence proved in
KT-AUXILIARY-UNIT-SELMER-BRIDGE: they kill the strict
groups with characters $\eta$, $\eta\chi^{-e}$ and
$\eta\chi^{-2e}$.  This section tests whether a
polarized generalized matrix algebra can turn that
information alone into exclusion of an irreducible
ordinary rank-three source.  The answer to that
\emph{formal implication} is negative.  The example
below is an abstract profinite group, not a number
field or an automorphic counterexample.

For the actual residual unitary action
$r_0=(\operatorname{diag}(\Psi,1,\Psi^{-1}),\eta)$,
$\Psi=\chi^e$, and
$\theta_{\rm Lie}(X)=-JX^tJ$, the root characters are
\[
\begin{array}{c|c}
\text{matrix line}&\text{global coefficient character}\\ \hline
E_{12}-E_{23}&\Psi\\
E_{12}+E_{23}&\eta\Psi\\
E_{21}-E_{32}&\Psi^{-1}\\
E_{21}+E_{32}&\eta\Psi^{-1}\\
E_{13}&\eta\Psi^2\\
E_{31}&\eta\Psi^{-2}.
\end{array}                                                \tag{EG.1}
\]
The negative rational line in this table is precisely
the original strict Selmer direction.  The positive
$\eta\Psi$ line is allowed by the upper ordinary flag;
it need not satisfy the strict local condition killed
by an auxiliary Bernoulli unit.  Even an additional
vanishing input for the rational positive-root group
does not remove these two remaining directions.

\begin{theorem}[KT-POLARIZED-EXTENSION-GRAPH-OBSTRUCTION]
Let $\ell>3$, let $e>0$ be even, and suppose that
the order $d$ of $\bar\chi^e$ in $\mathbf F_\ell^\times$
is greater than three.  There exist an abstract
profinite group $G$, characters $\chi_0,\eta$ on $G$,
and a continuous polarized representation
\[
 r:G\longrightarrow
          \mathrm{GL}_3(\mathbf Z_\ell)\rtimes\{\pm1\}
                                                               \tag{EG.2}
\]
with the following properties.  Put $\Psi=\chi_0^e$,
$G_E=\ker\eta$ and $k=\mathbf F_\ell$.

The reduction of $r$ is
$(\operatorname{diag}(\bar\Psi,1,\bar\Psi^{-1}),\eta)$,
and its characteristic-zero restriction to $G_E$ is
absolutely irreducible.  Among all one-dimensional
$k$-characters factoring through the finite quotient
constructed below, the only nonzero $H^1(G,k(\alpha))$
are
\[
 H^1(G,k(\eta\bar\Psi))\simeq k,
 \qquad H^1(G,k(\bar\Psi^{-1}))\simeq k.                  \tag{EG.3}
\]
In particular all of
\[
 H^1(G,k(\eta)),\quad
 H^1(G,k(\eta\bar\Psi^{-1})),\quad
 H^1(G,k(\eta\bar\Psi^{-2})),\quad
 H^1(G,k(\bar\Psi)),\quad
 H^1(G,k(\bar\Psi^2))                                    \tag{EG.4}
\]
vanish.  These five vanishings also hold with the corresponding
Teichm\"uller lifts of the characters and coefficients
$\mathbf Q_\ell/\mathbf Z_\ell$.
There is also a designated pair of
conjugate subgroups on which $r$ is upper triangular
in the prescribed order and the rational negative
class in (EG.3) restricts to zero.

Thus vanishing of these entire residual cohomology
groups, which is stronger than the corresponding
strict-group vanishing, is consistent with an
irreducible polarized representation and the indicated
ordinary flag.  No conclusion about actual Galois
groups or crystalline representations is asserted.
\end{theorem}

\begin{proof}
Choose a Teichm\"uller generator
$\zeta\in\mathbf Z_\ell^\times$ of order $\ell-1$
and put $q=\zeta^e$, of order $d>3$.
Let $F$ be the free pro-$\ell$ group on generators
$x,y$, and let
$H=\langle h\rangle\times\langle c\rangle
\simeq C_{\ell-1}\times C_2$ act by
\[
 hxh^{-1}=x^q,\quad hyh^{-1}=y^{q^{-1}},\qquad
 cxc=x^{-1},\quad cyc=y.                                 \tag{EG.5}
\]
Powers by $\mathbf Z_\ell$-units are defined in the
procyclic subgroups generated by $x,y$.
The indicated substitutions are automorphisms: their
inverse substitutions use $q^{-1},q$.  They commute,
and the substitution for $h$ has order dividing
$\ell-1$.  Hence $G=F\rtimes H$ is well defined.
Set $\chi_0(h)=\zeta$, $\chi_0(c)=-1$,
$\chi_0(F)=1$, and $\eta(c)=-1$, $\eta(h)=\eta(F)=1$.
Evenness of $e$ gives $\Psi(c)=1$.

Set
\[
 X=E_{12}+E_{23},\quad Y=E_{21}-E_{32},\quad
 T=\operatorname{diag}(q,1,q^{-1}).                       \tag{EG.6}
\]
Define
\[
 r(h)=(T,1),\quad r(c)=(1,-1),\quad
 r(x)=(\exp(\ell X),1),\quad
 r(y)=(\exp(\ell Y),1).                                  \tag{EG.7}
\]
The exponentials terminate in degree two and are
integral because two is a unit.  Their images lie in
the pro-$\ell$ group $1+\ell M_3(\mathbf Z_\ell)$,
so freeness gives the required continuous map on $F$.
Conjugation by $T$ scales $X,Y$ by $q,q^{-1}$.
Moreover $\theta_{\rm Lie}(X)=-X$ and
$\theta_{\rm Lie}(Y)=Y$.  Thus (EG.7) respects every
relation (EG.5) and defines (EG.2), including its
polarization and residual normalization.

For absolute irreducibility it suffices to work over
$\overline{\mathbf Q}_\ell$.  Since the eigenvalues
$q,1,q^{-1}$ of $T$ are distinct, a subspace invariant
under $T$ is spanned by a subset of the coordinate
vectors.  Invariance under $\exp(\ell X)$ and
$\exp(\ell Y)$ then forces every nonempty such subset
to contain all three vectors.  Equivalently, these
two logarithms generate $\mathfrak{sl}_3$:
\[
 [X,Y]=\operatorname{diag}(1,-2,1)=Z,\quad
 [Z,X]=3(E_{12}-E_{23}),\quad
 [Z,Y]=-3(E_{21}+E_{32}).                                 \tag{EG.8}
\]
These identities recover the four simple matrix
units, their brackets the two long units and the
two diagonal directions.

For (EG.3), $F$ acts trivially on each residual
coefficient module, and $\ell\nmid|H|$.
Inflation--restriction and averaging therefore give
\[
 H^1(G,k(\alpha))
 \simeq\operatorname{Hom}_H(F^{\rm ab}/\ell,
                                      k(\alpha)).          \tag{EG.9}
\]
The two generators of $F^{\rm ab}/\ell$ have exactly
the characters $\eta\bar\Psi$ and
$\bar\Psi^{-1}$ under $H$ by (EG.5).
This proves (EG.3).  None of the characters in
(EG.4) equals either of them when $d>3$.
For the divisible-coefficient assertion, each named
character is nontrivial modulo $\ell$, so its global
$H^0$ is zero.  The multiplication-by-$\ell$ exact
sequence identifies the $\ell$-torsion in its
divisible-coefficient $H^1$ with the corresponding
zero group in (EG.4).  Continuous cocycles from a
compact group into a discrete torsion module have
finite image, so this $H^1$ is $\ell$-primary torsion
and must itself be zero.

Finally let $D$ be the closed subgroup generated by
$h,x$.  Then $cDc=D$, so these give a designated
conjugate pair.  Its image under $r$ is upper
triangular with diagonal order $(\Psi,1,\Psi^{-1})$.
The cocycle corresponding to the second generator
of $F^{\rm ab}/\ell$ is zero on $D$, proving the
claimed strictness.  This is a group-theoretic local
analogue only; the theorem does not replace $D$ by
an actual local Galois group.
\end{proof}

The obstruction is the mixed pair of arrows
$E_{12}+E_{23}$ and $E_{21}-E_{32}$.  It connects all
three residual constituents.  Their commutators create
the missing odd negative and long-root directions as
higher products, as (EG.8) shows.  Vanishing of a
primitive extension quotient $A_{ij}/A'_{ij}$ can
therefore leave $A_{ij}=A'_{ij}$ nonzero.  Applying
Nakayama to the whole off-diagonal module without
first proving that the path products lie in its
coefficient maximal ideal would lose this distinction.

For the arithmetic problem, an additional argument
must control those mixed products, or relate an
integral extension module to an independently bounded
arithmetic quantity.  Requiring the remaining rational
negative strict group itself to vanish would be the
original Selmer hypothesis.  The three field units,
even together with rational positive-root vanishing,
do not establish that assertion by the formal GMA
argument tested here.
